
\documentclass{article}
\usepackage{colortbl}
\usepackage{makecell}
\usepackage{multirow}
\usepackage{supertabular}

\begin{document}

\newcounter{utterance}

\centering \large Interaction Transcript for game `codenames', experiment `association\_difficult', episode 1 with reflexion{-}Qwen3.5{-}4B.
\vspace{24pt}

{ \footnotesize  \setcounter{utterance}{1}
\setlength{\tabcolsep}{0pt}
\begin{supertabular}{c@{$\;$}|p{.15\linewidth}@{}p{.15\linewidth}p{.15\linewidth}p{.15\linewidth}p{.15\linewidth}p{.15\linewidth}}
    \# & $\;$A & \multicolumn{4}{c}{Game Master} & $\;\:$B\\
    \hline

    \theutterance \stepcounter{utterance}  
    & & & \multicolumn{2}{p{0.3\linewidth}}{
        \cellcolor[rgb]{0.9,0.9,0.9}{
            \makecell[{{p{\linewidth}}}]{
                \texttt{\tiny{[GM$|$GM]}}
                {'hidden': {'team': ['spider', 'admiral', 'back', 'xylophone'], 'innocent': ['ear', 'organ', 'glacier', 'butter', 'lungs', 'firefighter', 'sweatshirt'], 'opponent': [], 'assassin': ['banker']}, 'revealed': {'team': {'team': [], 'innocent': [], 'opponent': [], 'assassin': []}, 'opponent': {'team': [], 'innocent': [], 'opponent': [], 'assassin': []}}}
            }
        }
    }
    & & \\ \\

    \theutterance \stepcounter{utterance}  
    & & \multicolumn{4}{p{0.6\linewidth}}{
        \cellcolor[rgb]{0.9,0.9,0.9}{
            \makecell[{{p{\linewidth}}}]{
                \texttt{\tiny{[P1$\langle$GM]}}
                \texttt{You have attempted this same task before and did not succeed. These are the plans you wrote after those attempts.} \\
\texttt{Trial \#1: New plan: The last attempt was aborted because the response format was incorrect; specifically, the TARGETS line was missing the required "TARGETS: " prefix. Next time, I will strictly ensure the output follows the exact two{-}line format: "CLUE: <WORD>" followed by a newline and "TARGETS: <WORD1>, <WORD2>". I will also verify that the clue word is semantically related to the targets and does not appear in the distractor or kill word lists before sending.} \\
\\ 
\texttt{Trial \#2: New plan: I failed because I guessed "admiral" using the clue "ARMY," which is semantically related but too broad and likely led to the opponent guessing it as a distractor or kill word, or simply wasted a turn without securing the win. The strategy should prioritize clues that are unique to the remaining team words and less likely to overlap with opponent words or distractors. For the remaining words (spider, back, xylophone), I should use specific associations like "WEB" for spider, "POST" or "REST" for back, or "KEY" or "INSTRUMENT" for xylophone, ensuring the clue is distinct enough to prevent the opponent from guessing them incorrectly. I will also avoid using words that might be confused with the opponent's remaining word "daisy" or the kill word "banker."} \\
\\ 
\texttt{Trial \#3: New plan: I failed because I only provided one target ("spider") when I have four team words ("spider, admiral, back, xylophone"). The game rules state that providing more related words allows for a quicker win, but I must ensure the clue is specific enough to avoid the opponent guessing my words as distractors or kill words. Next time, I will select a single{-}word clue that uniquely relates to at least two of my remaining words (e.g., "BACK" for "back" and "admiral", or "KEY" for "xylophone" and potentially others if applicable) and list all related targets in the TARGETS line, ensuring the clue does not match any distractor or kill word.} \\
\\ 
\\ 
\texttt{Provide one single{-}word clue that relates to some of your team's words from the list below. You can choose to how many and to which words your clue relates to.} \\
\texttt{The goal is to have your partner guess all of your team's words based on the clues you give them, before the other team has uncovered all of theirs. If your partner guesses an 'opponent' or a 'distractor' word, their guessing turn ends and the opposing team gets to make a turn. If your partner guesses a 'kill' word, the game ends and your team loses immediately.} \\
\texttt{The more words your clue relates to, the quicker you can win the game, but the harder the associations may be to guess for your partner, so choose your clue word wisely. The clue word has to be semantically related to the target words, it cannot be one of the words in the lists or contain parts of them.} \\
\\ 
\texttt{Always give your single{-}word clue and your comma{-}separated list of related target words in the following format and make your answers as short as possible, never include any other text than is required in this form:} \\
\\ 
\texttt{CLUE: <WORD>} \\
\texttt{TARGETS: <TARGETS>} \\
\\ 
\texttt{Your team words are: spider, admiral, back, xylophone.} \\
\texttt{Your opponent's team words are: .} \\
\texttt{Distractor words are: ear, organ, glacier, butter, lungs, firefighter, sweatshirt.} \\
\texttt{The kill words are: banker.} \\
            }
        }
    }
    & \\ \\

    \theutterance \stepcounter{utterance}  
    & \multicolumn{4}{p{0.6\linewidth}}{
        \cellcolor[rgb]{0.9,0.9,0.9}{
            \makecell[{{p{\linewidth}}}]{
                \texttt{\tiny{[P1$\rangle$GM]}}
                \texttt{CLUE: WEB} \\
\texttt{TARGETS: spider, admiral} \\
            }
        }
    }
    & & \\ \\

    \theutterance \stepcounter{utterance}  
    & & & \multicolumn{2}{p{0.3\linewidth}}{
        \cellcolor[rgb]{0.9,0.9,0.9}{
            \makecell[{{p{\linewidth}}}]{
                \texttt{\tiny{[GM$|$GM]}}
                \texttt{web} \\
            }
        }
    }
    & & \\ \\

    \theutterance \stepcounter{utterance}  
    & & & \multicolumn{2}{p{0.3\linewidth}}{
        \cellcolor[rgb]{0.9,0.9,0.9}{
            \makecell[{{p{\linewidth}}}]{
                \texttt{\tiny{[GM$|$GM]}}
                ['spider', 'admiral']
            }
        }
    }
    & & \\ \\

    \theutterance \stepcounter{utterance}  
    & & & \multicolumn{2}{p{0.3\linewidth}}{
        \cellcolor[rgb]{0.9,0.9,0.9}{
            \makecell[{{p{\linewidth}}}]{
                \texttt{\tiny{[GM$|$GM]}}
                {'word': 'spider', 'assignment': 'team'}
            }
        }
    }
    & & \\ \\

    \theutterance \stepcounter{utterance}  
    & & & \multicolumn{2}{p{0.3\linewidth}}{
        \cellcolor[rgb]{0.9,0.9,0.9}{
            \makecell[{{p{\linewidth}}}]{
                \texttt{\tiny{[GM$|$GM]}}
                {'word': 'admiral', 'assignment': 'team'}
            }
        }
    }
    & & \\ \\

    \theutterance \stepcounter{utterance}  
    & & & \multicolumn{2}{p{0.3\linewidth}}{
        \cellcolor[rgb]{0.9,0.9,0.9}{
            \makecell[{{p{\linewidth}}}]{
                \texttt{\tiny{[GM$|$GM]}}
                \texttt{opponent has won} \\
            }
        }
    }
    & & \\ \\

    \theutterance \stepcounter{utterance}  
    & \multicolumn{4}{p{0.6\linewidth}}{
        \cellcolor[rgb]{0.95,0.95,0.95}{
            \makecell[{{p{\linewidth}}}]{
                \texttt{\tiny{[P2$\rightarrow$P1]}}
                \texttt{Post{-}episode Reflexion (not sent to GM):} \\
\texttt{New plan: I failed because I only listed one target ("spider") when I have four team words, missing the opportunity to reveal multiple words with a single clue. Next time, I will select a single{-}word clue that uniquely relates to at least two of my remaining words (e.g., "BACK" for "back" and "admiral", or "KEY" for "xylophone" and "xylophone" if applicable) and list all related targets in the TARGETS line, ensuring the clue does not match any distractor or kill word. I will prioritize clues that cover the most remaining words to maximize the speed of the game.} \\
            }
        }
    }
    & & \\ \\

\end{supertabular}
}

\end{document}
